% Generated draft: requires mathematical and source review.
\documentclass[11pt]{article}
\usepackage[margin=25mm]{geometry}
\usepackage{fontspec}
\setmainfont{Latin Modern Roman}
\usepackage{amsmath,amssymb,mathtools,unicode-math}
\setmathfont{Latin Modern Math}
\usepackage{enumitem,needspace,xurl,hyperref}
\setlength{\parindent}{0pt}
\setlength{\parskip}{5pt}
\setlength{\emergencystretch}{3em}
\newcommand{\N}{\mathbb{N}}
\newcommand{\Z}{\mathbb{Z}}
\newcommand{\Q}{\mathbb{Q}}
\newcommand{\R}{\mathbb{R}}
\newcommand{\C}{\mathbb{C}}
\newcommand{\F}{\mathbb{F}}
\newcommand{\A}{\mathbb{A}}
\newcommand{\PP}{\mathbb{P}}
\newcommand{\E}{\mathbb{E}}
\newcommand{\eps}{\varepsilon}
\newcommand{\dd}{\,\mathrm{d}}
\begin{document}
\section*{UnsolvedMath ID: 701266}
\textbf{Source problem code:} catalogue entry
\subsection{Definitions and mathematical statement}
The principal target is part (b) of Kourovka Notebook Problem 16.74, proposed by V. V. Nekrashevich in 2006.

A group \(G\) is \emph{amenable} if there exists a finitely additive probability measure \(\mu\colon\mathcal P(G)\to[0,1]\) such that \(\mu(gA)=\mu(A)\) for every \(g\in G\) and every \(A\subseteq G\).

Part (b) asks whether every group generated by automata of polynomial growth, in the sense of S. Sidki's cited 2004 paper, is amenable. The source does not reproduce Sidki's definition. A conventional polynomial-activity interpretation can be made precise as follows. Let \(X\) be a finite nonempty alphabet, and let \(X^*\) be the rooted tree of finite words over \(X\). For a rooted-tree automorphism \(g\) and a word \(v\), its section \(g|_v\) is characterized by
\[
 g(vw)=g(v)\,g|_v(w)\qquad(w\in X^*).
\]
An automorphism is finite-state if its set of sections is finite. Define
\[
 a_g(n)=\bigl|\{v\in X^n:g|_v\ne\mathrm{id}\}\bigr|.
\]
It has polynomial activity if there exist \(C>0\) and \(d\in\Z_{\ge0}\) such that
\[
 a_g(n)\le C(n+1)^d\qquad\text{for every }n\in\Z_{\ge0}.
\]
Under this interpretation, the question is whether, for every such alphabet \(X\) and every set \(S\subseteq\operatorname{Aut}(X^*)\) of finite-state automorphisms of polynomial activity, the group \(\langle S\rangle\) is amenable. The constants and degree may depend on the generator. This is an explicit interpretation requiring comparison with Sidki's precise convention; the supplied source does not specify whether the generating automata must form a finite family or share a uniform polynomial bound.

The record also retains archived part (a). It describes a purported group \(G=\langle\alpha,\beta\rangle\) of permutations of \(\Z\), with formulas
\[
 \alpha(n)=n+1,\qquad \beta(0)=0,\qquad
 \beta(2^k m)=2^k(m+2),
\]
and asks whether \(G\) is amenable. The supplied qualification is that \(m\) is odd and \(n\) is a positive integer. No range for \(k\) is supplied, and the positive-integer restriction on \(n\) does not specify \(\alpha\) on all of \(\Z\). Consequently, these formulas are preserved without treating them as a complete definition of the claimed permutation group. The supplied archive reports an affirmative answer to part (a), attributed to G. Amir, O. Angel and B. Vir\'ag; that report is not independently verified here.
\subsection{Short English statement}
Part (b) asks whether every group generated by polynomial-growth automata in Sidki's sense is amenable. A conventional polynomial-activity formulation is provided, with its correspondence to the cited definition left for review. The supplied record labels the problem partially solved and reports an affirmative solution to archived part (a), whose displayed generator definitions contain unresolved domain qualifications.
\subsection{Sources}
\begin{itemize}
\item Local UnsolvedMath registry, permanent ID \(701266\), exact source code \texttt{catalogue entry}, title \emph{Kourovka Notebook Problem 16.74}; supplied complete source record, upstream revision \texttt{d371ee917ac9647fbb06cd6fc529deec6f3dda82}.
\item \emph{The Kourovka Notebook (Unsolved Problems in Group Theory)}, 16th issue (2006), Problem 16.74, proposed by V. V. Nekrashevich. The supplied record identifies an update in the 21st issue (2026), \texttt{arXiv:1401.0300v48}. The collection is attributed to editors E. I. Khukhro and V. D. Mazurov.
\item S. Sidki, \emph{Geom. Dedicata} \textbf{108} (2004), 193--204. Supplied reference for the meaning of automata of polynomial growth; no article title is supplied.
\item G. Amir, O. Angel and B. Vir\'ag, \emph{J. Eur. Math. Soc.} \textbf{15}, no. 3 (2013), 705--730. The supplied Archive of Solved Problems attributes an affirmative answer to part (a) to this reference; no article title is supplied.
\item Supplied status review dated \texttt{2026-10-16}: canonical status \texttt{partially\_solved}, evidence label \texttt{published}; part (a) is reported solved in the archive, with the remaining part in the main section. This is supplied dated curation, not an independently verified literature assessment. The cited publications and archive have not been consulted for this draft.
\end{itemize}
\end{document}
